% Options for packages loaded elsewhere
\PassOptionsToPackage{unicode}{hyperref}
\PassOptionsToPackage{hyphens}{url}
\PassOptionsToPackage{dvipsnames,svgnames,x11names}{xcolor}
%
\documentclass[
  10pt,
]{article}
\usepackage{amsmath,amssymb}
\usepackage{iftex}
\ifPDFTeX
  \usepackage[T1]{fontenc}
  \usepackage[utf8]{inputenc}
  \usepackage{textcomp} % provide euro and other symbols
\else % if luatex or xetex
  \usepackage{unicode-math} % this also loads fontspec
  \defaultfontfeatures{Scale=MatchLowercase}
  \defaultfontfeatures[\rmfamily]{Ligatures=TeX,Scale=1}
\fi
\usepackage{lmodern}
\ifPDFTeX\else
  % xetex/luatex font selection
\fi
% Use upquote if available, for straight quotes in verbatim environments
\IfFileExists{upquote.sty}{\usepackage{upquote}}{}
\IfFileExists{microtype.sty}{% use microtype if available
  \usepackage[]{microtype}
  \UseMicrotypeSet[protrusion]{basicmath} % disable protrusion for tt fonts
}{}
\makeatletter
\@ifundefined{KOMAClassName}{% if non-KOMA class
  \IfFileExists{parskip.sty}{%
    \usepackage{parskip}
  }{% else
    \setlength{\parindent}{0pt}
    \setlength{\parskip}{6pt plus 2pt minus 1pt}}
}{% if KOMA class
  \KOMAoptions{parskip=half}}
\makeatother
\usepackage{xcolor}
\usepackage[margin=24mm]{geometry}
\usepackage{longtable,booktabs,array}
\usepackage{calc} % for calculating minipage widths
% Correct order of tables after \paragraph or \subparagraph
\usepackage{etoolbox}
\makeatletter
\patchcmd\longtable{\par}{\if@noskipsec\mbox{}\fi\par}{}{}
\makeatother
% Allow footnotes in longtable head/foot
\IfFileExists{footnotehyper.sty}{\usepackage{footnotehyper}}{\usepackage{footnote}}
\makesavenoteenv{longtable}
\usepackage{graphicx}
\makeatletter
\def\maxwidth{\ifdim\Gin@nat@width>\linewidth\linewidth\else\Gin@nat@width\fi}
\def\maxheight{\ifdim\Gin@nat@height>\textheight\textheight\else\Gin@nat@height\fi}
\makeatother
% Scale images if necessary, so that they will not overflow the page
% margins by default, and it is still possible to overwrite the defaults
% using explicit options in \includegraphics[width, height, ...]{}
\setkeys{Gin}{width=\maxwidth,height=\maxheight,keepaspectratio}
% Set default figure placement to htbp
\makeatletter
\def\fps@figure{htbp}
\makeatother
\setlength{\emergencystretch}{3em} % prevent overfull lines
\providecommand{\tightlist}{%
  \setlength{\itemsep}{0pt}\setlength{\parskip}{0pt}}
\setcounter{secnumdepth}{-\maxdimen}
\ifLuaTeX
\usepackage[bidi=basic]{babel}
\else
\usepackage[bidi=default]{babel}
\fi
\babelprovide[main,import]{english}
% get rid of language-specific shorthands (see #6817):
\let\LanguageShortHands\languageshorthands
\def\languageshorthands#1{}
\usepackage{microtype}
\usepackage{float}
\usepackage{xurl}
\floatplacement{figure}{htbp}
\setlength{\emergencystretch}{3em}
\ifLuaTeX
  \usepackage{selnolig}  % disable illegal ligatures
\fi
\usepackage{bookmark}
\IfFileExists{xurl.sty}{\usepackage{xurl}}{} % add URL line breaks if available
\urlstyle{same}
\hypersetup{
  pdftitle={Source-bound knowledge consolidation and recognition sensitivity in channeling discourse},
  pdfauthor={Maintainer/authorship: TBD --- AI-assisted technical report},
  pdflang={en},
  colorlinks=true,
  linkcolor={Maroon},
  filecolor={Maroon},
  citecolor={Blue},
  urlcolor={Blue},
  pdfcreator={LaTeX via pandoc}}

\title{Source-bound knowledge consolidation and recognition sensitivity
in channeling discourse}
\usepackage{etoolbox}
\makeatletter
\providecommand{\subtitle}[1]{% add subtitle to \maketitle
  \apptocmd{\@title}{\par {\large #1 \par}}{}{}
}
\makeatother
\subtitle{An additive, reproducible technical supplement to Channels ET}
\author{Maintainer/authorship: TBD --- AI-assisted technical report}
\date{8 November 2026}

\begin{document}
\maketitle

\section{Abstract}\label{abstract}

We examine how a small, source-bound knowledge bank changes when texts
are represented lexically, embedded, consolidated into qualified answer
cards, and transcribed by a different recognition pipeline. The frozen
frame contains \textbf{25 distinct source editions, 47,280 whitespace
words and 6 medium contexts}. Close reading covers 15,795 unique body
words in 14 sources; it does not cover the entire modeled frame. The
public bank contains 13 question labels, 25 answer cards and 11 scoped
comparisons. A typed graph has 100 nodes and 135 edges.

Lexical and neural neighborhoods agree only partially: mean top-five
Jaccard is 0.340. Same-medium pairs are descriptively more similar than
cross-medium pairs sharing a comparison label; format, editing and
speaker context remain confounded. A four-query source-retrieval
diagnostic exposes a mathematical coverage bias in naive full-list
reciprocal-rank fusion (RRF). A separately labelled post-hoc repair does
not establish a general retrieval winner. A paired eight-minute
speech-recognition pilot runs NVIDIA Parakeet on two CPU workers and
compares its output with Whisper tiny.en. Parakeet takes 62.40 seconds
versus 12.48; disagreement is measured without treating either output as
ground truth. A separate completed four-recording run covers 178.99
audio minutes and yields 26,319 recognized words; downstream sensitivity
preserves older editions and requires re-reading affected cards. This
study updates no language-model weights and does not reproduce Google's
Sleep training framework. It supports auditable navigation, source
criticism and experiment design, rather than verification of the
identities or events claimed in channeling discourse.

\textbf{Status.} This is an AI-assisted technical report, not a
submitted or peer-reviewed paper. Maintainer authorship, institutional
affiliations and human annotation validation remain TBD. Earlier
corpora, papers and source editions remain unchanged.

\clearpage
\tableofcontents
\clearpage

\section{1. Questions and study
boundaries}\label{questions-and-study-boundaries}

The immediate research questions are methodological. Do lexical and
neural representations organize the same editions similarly? Does
preserving the conditions of an answer improve useful retrieval, or
merely increase index coverage and verbosity? Can two statements be
compared without erasing their time, audience, referent or modal force?
Does a larger contemporary speech-recognition model change downstream
representations enough to require new annotation?

The unit of analysis is a \textbf{documentary source edition}, not an
extraterrestrial civilization, a verified spiritual speaker, a unique
audience member or a complete YouTube channel. A medium is a recorded
presentation or attribution context. A comparison label is an analyst
grouping of printed identities; broad civilization groups and named
personalities occupy different conceptual levels. The frame contains 13
comparison labels and 17 exact printed labels. Neither shared wording
nor a graph edge authenticates a claimed identity.

The 25-edition sample joins a previous 20-document written pilot with
five Council-related texts. It is a purposive, accessible-source frame.
It is not the separate, previously completed 228-unit publisher fit, a
random sample of channeling, or an inventory of all videos. No new
source media were acquired to assemble this particular frame. Existing
retained audio was used for the recognition experiments. Duplicate
checks removed no editions in this selection; this does not establish
independence of their underlying ideas, publishing workflows or
audiences.

\section{2. Literature and the meaning of
consolidation}\label{literature-and-the-meaning-of-consolidation}

\subsection{2.1 The identifiable Google
paper}\label{the-identifiable-google-paper}

The strongest match to the requested recent Google work is
\textbf{Behrouz, Hashemi, Javanmard and Mirrokni, ``Language Models Need
Sleep: Learning to Self-Modify and Consolidate Memories,''
arXiv:2606.03979v2} {[}R1{]}. Version 1 appeared on 2 June 2026, version
2 on 10 July 2026; the metadata also notes an earlier OpenReview version
from September 2025. Google's Research blog describes it on \textbf{27
October 2026} {[}R2{]}. The identification is source-grounded but
remains a likely match because no paper link was supplied by the
requester.

The authors' framework performs parameter updates. A fast wake memory
operates alongside slower memory modules. It expands or activates
low-rank mixture-of-experts memory while preserving older parameters,
transfers information through generalized on-policy distillation and
reinforcement-learning imitation, and resets fast modules after
transfer. Its dreaming procedure activates random additional experts,
produces synthetic curricula, selects examples using gradients, and
trains adaptations. Cross-field connection in this setting is a proposed
computational benefit of trained memory recombination; it is not
evidence that two discourse traditions share an external referent.

The reported benchmarks require their original denominators. The
ten-million-token BABILong result uses task-finetuned small models, and
Appendix B.2 reports substantial degradation without that fine-tuning.
The filtered ARC experiment uses 11 training tasks and eight held-out
tasks, reporting successful-dream fractions from five dreams per test
task. Mathematical results use the stated average-at-16 setting.
Appendix B.5 reports same-step supervised fine-tuning as four times more
efficient, while selected time-to-target comparisons have a different
ordering. None of those values forecasts performance on this corpus or
CPU environment. The inspected paper and blog did not establish an
executable official author release; similarly named third-party
repositories are not treated as reproductions.

Our implementation is an \textbf{external source-card and index
consolidation experiment}. Sources remain immutable, cards are explicit
annotations, and retrieval is measured separately. There are no
trainable memory experts, imitation training, synthetic-gradient
selection, new model weights, fine-tuning or biological sleep
measurements. The analogy is limited to preparing a persistent,
compressed, auditable representation before later queries.

\subsection{2.2 Adjacent primary
methods}\label{adjacent-primary-methods}

Sleep-time Compute {[}R3{]} is a distinct Letta/UC Berkeley line of work
about offline context preparation; it must not be attributed to Google.
Nested Learning {[}R4{]} supplies architectural context for nested
update frequencies, not proof that storing a card implements continual
learning. GraphRAG {[}R5{]} motivates document-to-entity-to-community
summaries for global questions. HippoRAG {[}R6{]} links passage
retrieval, OpenIE and personalized PageRank; HippoRAG 2 {[}R7{]}
broadens factual, associative and sensemaking retrieval. RAPTOR {[}R8{]}
motivates hierarchical retrieval over summaries. These are design
references; the present typed graph is not an executed GraphRAG,
HippoRAG or RAPTOR benchmark.

MMTEB {[}R9{]} motivates task-specific, multilingual embedding
evaluation rather than claiming a universally best encoder.
Sentence-BERT {[}R10{]} supplies the foundational rationale for
sentence-level similarity. The local encoder is a small, pinned MiniLM
baseline, not a current universal state-of-the-art claim. Work on
contradictions in text {[}R11{]} makes referent, scope and modality
central: negation or cosine alone cannot decide contradiction. Whisper
{[}R12{]} provides speech-recognition background, but the inspected
abstract is not evidence for a specific observed error mechanism.
Datasheets {[}R13{]} and W3C PROV-O {[}R14{]} motivate explicit
collection, processing and agent provenance. The JSON records borrow
that conceptual separation without claiming formal RDF conformance.

Access depth is recorded. HippoRAG 2, RAPTOR, Sentence-BERT, Whisper and
Datasheets were inspected at abstract level in this literature pass. The
report does not imply full-paper validation of their detailed methods.
The literature register contains the exact versions, retrieval dates,
inspected sections and source hashes for all 14 primary sources.

\section{3. Data, provenance and
annotation}\label{data-provenance-and-annotation}

\subsection{3.1 Coverage is not a single
denominator}\label{coverage-is-not-a-single-denominator}

\begin{longtable}[]{@{}
  >{\raggedright\arraybackslash}p{(\columnwidth - 2\tabcolsep) * \real{0.5000}}
  >{\raggedright\arraybackslash}p{(\columnwidth - 2\tabcolsep) * \real{0.5000}}@{}}
\toprule\noalign{}
\begin{minipage}[b]{\linewidth}\raggedright
Operation
\end{minipage} & \begin{minipage}[b]{\linewidth}\raggedright
Units actually covered
\end{minipage} \\
\midrule\noalign{}
\endhead
\bottomrule\noalign{}
\endlastfoot
Representation analysis & 25 editions / 47,280 body words / 6 medium
contexts \\
Close reading & 14 editions / 15,795 unique body words \\
Complete retained source bodies read & 10 \\
Additional documentary question text & 65 words outside the body
denominator \\
Question bank & 8 analyst-authored + 5 paraphrased printed questions \\
Knowledge cards and comparisons & 25 cards / 11 comparison records \\
Typed graph & 100 nodes / 135 edges \\
Fixed retrieval diagnostic & 4 source-informed questions; one known
support per question \\
Paired ASR pilot & 16 windows / 8 recordings / 8 unique audio minutes \\
\end{longtable}

The frozen frame's medium-level distribution is:

\begin{longtable}[]{@{}
  >{\raggedright\arraybackslash}p{(\columnwidth - 6\tabcolsep) * \real{0.2500}}
  >{\raggedright\arraybackslash}p{(\columnwidth - 6\tabcolsep) * \real{0.2500}}
  >{\raggedright\arraybackslash}p{(\columnwidth - 6\tabcolsep) * \real{0.2500}}
  >{\raggedright\arraybackslash}p{(\columnwidth - 6\tabcolsep) * \real{0.2500}}@{}}
\toprule\noalign{}
\begin{minipage}[b]{\linewidth}\raggedright
Medium attribution context
\end{minipage} & \begin{minipage}[b]{\linewidth}\raggedright
Editions
\end{minipage} & \begin{minipage}[b]{\linewidth}\raggedright
Words
\end{minipage} & \begin{minipage}[b]{\linewidth}\raggedright
Machine transcript editions
\end{minipage} \\
\midrule\noalign{}
\endhead
\bottomrule\noalign{}
\endlastfoot
Daniel Scranton & 8 & 4,394 & 0 \\
Lyssa Royal Holt & 1 & 2,671 & 0 \\
Natalie Glasson & 8 & 6,927 & 0 \\
Pamela Aaralyn & 1 & 16,931 & 1 \\
Pamela Rose Kribbe & 4 & 7,324 & 0 \\
Tracy Kumbera & 3 & 9,033 & 3 \\
\end{longtable}

This table counts the retained documentary editions, not a medium's
complete output. The source mix is uneven; the machine-transcript column
also exposes a major text-origin confound. The five documentary question
labels paraphrase printed questions attributed to the Lyssa Royal
Holt/Sasha source; the other eight are analyst-written after reading
their answers. They are navigation labels, not independently observed
matching questions asked to multiple mediums. The analyst questions are
excluded from the experimental card index. Automated full-text
processing does not count as close reading. No human listening
validation or human consensus annotation is claimed.

Source records retain edition identifiers, retrieval context, URL and
cryptographic bindings. Normalized text locators use NFKC-normalized,
stripped Python strings with end-exclusive character offsets. Source
wording is retained privately; public artifacts expose bounded
paraphrases, identifiers, evidence locators, summaries and plots.
Existing Daniel attribution aliases are resolved to the same person
where documented. Pamela Aaralyn and Pamela Rose Kribbe remain distinct.
Recognition outputs retain unknown speaker roles where speaker identity
is not established.

\subsection{3.2 Cards preserve
conditions}\label{cards-preserve-conditions}

A card separates a question label, a concise answer, a scope or
qualifier, and support locators. A relation compares propositions, not
entities. The annotation record contains:

\begin{longtable}[]{@{}ll@{}}
\toprule\noalign{}
Annotation relation & Pairs \\
\midrule\noalign{}
\endhead
\bottomrule\noalign{}
\endlastfoot
compatible & 4 \\
conditional tension & 1 \\
different focus & 3 \\
explicit contradiction & 0 \\
shared proposition & 1 \\
unresolved & 2 \\
\end{longtable}

Examples show why qualification matters. A source's advice to trust the
Universe or Creator does not license the broader claim that every
interpersonal relationship is safe. Comparing it with a boundary-setting
passage yields a conditional tension only if an analyst expands the
referent to draining relationships; it is not an explicit contradiction
in the retained texts. Optional invocation of an Andromedan framing can
coexist with a nonhierarchical Jeshua account. Voluntary recentering can
coexist with rejection of a compulsory universal formula. A contact
passage's near-term discussion concerns isolated contact and does not
establish a definite start date for universal contact. These are bounded
interpretive comparisons, not endorsements of the source claims.

Independent review corrected qualifiers and locators before acceptance.
The graph contains typed source, annotation, question, topic and
methodological-hypothesis links. Its printed-label nodes are
medium-scoped: Mother Mary and Mary Magdalene each remain separate for
Kribbe and Glasson. Consequently, 19 graph label nodes coexist with 17
distinct literal label strings in the analysis. Its 135 edges are not
135 independent corroborations, causal relationships or verified
spiritual interactions. A node's degree partly reflects annotation
choices and source coverage.

\section{4. Representations, topics and source
effects}\label{representations-topics-and-source-effects}

\subsection{4.1 Fixed processing
contract}\label{fixed-processing-contract}

Complete exact medium-name and printed-label phrases are masked before
the main similarity analysis. This reduces a direct name cue but leaves
aliases, vocabulary, repeated introductions, editing practices and
source format. TF-IDF uses unigram and bigram features, sublinear term
frequency, smoothed inverse-document frequency, L2 normalization,
minimum document frequency one and a 12,000-feature ceiling. Unicode
letters and apostrophes define tokens; the stop list retains ``no,''
``not,'' ``nor'' and ``never.'' Standalone numerals are not lexical
features. Numerical claims therefore require a separate source-level
audit.

The semantic encoder is \nolinkurl{sentence-transformers/all-MiniLM-L6-v2},
pinned to revision \nolinkurl{1110a243fdf4706b3f48f1d95db1a4f5529b4d41},
running through cached ONNX on one CPU thread. Documents are encoded as
nonoverlapping 240-content-wordpiece chunks plus special tokens.
Attention-mask mean pooling includes those special tokens, each chunk is
L2 normalized, and the document mean is weighted by content-wordpiece
count before final L2 normalization. All content chunks are processed;
this is not head-only truncation. Equal-chunk weighting is a sensitivity
condition.

A second condition uses a centered 300-analysis-token span selected
before name masking. All 25 editions qualify. It reuses the full-corpus
lexical vocabulary and IDF rather than refitting. This is a fixed-span
control, not a random excerpt sample. Three NMF factors and three
K-means clusters are fixed exploratory displays, not statistically
selected doctrine counts. K-means uses ten initializations for seeds 11,
23 and 47.

\subsection{4.2 Similarity is partly a source-context
measurement}\label{similarity-is-partly-a-source-context-measurement}

There are 300 unordered pairs, and pairs sharing an edition are
dependent. Their descriptive means are:

\begin{longtable}[]{@{}
  >{\raggedright\arraybackslash}p{(\columnwidth - 10\tabcolsep) * \real{0.1667}}
  >{\raggedright\arraybackslash}p{(\columnwidth - 10\tabcolsep) * \real{0.1667}}
  >{\raggedright\arraybackslash}p{(\columnwidth - 10\tabcolsep) * \real{0.1667}}
  >{\raggedright\arraybackslash}p{(\columnwidth - 10\tabcolsep) * \real{0.1667}}
  >{\raggedright\arraybackslash}p{(\columnwidth - 10\tabcolsep) * \real{0.1667}}
  >{\raggedright\arraybackslash}p{(\columnwidth - 10\tabcolsep) * \real{0.1667}}@{}}
\toprule\noalign{}
\begin{minipage}[b]{\linewidth}\raggedright
Pair stratum
\end{minipage} & \begin{minipage}[b]{\linewidth}\raggedright
Pairs
\end{minipage} & \begin{minipage}[b]{\linewidth}\raggedright
TF-IDF
\end{minipage} & \begin{minipage}[b]{\linewidth}\raggedright
MiniLM
\end{minipage} & \begin{minipage}[b]{\linewidth}\raggedright
Span TF-IDF
\end{minipage} & \begin{minipage}[b]{\linewidth}\raggedright
Span MiniLM
\end{minipage} \\
\midrule\noalign{}
\endhead
\bottomrule\noalign{}
\endlastfoot
Different medium, different label & 222 & 0.054 & 0.545 & 0.023 &
0.380 \\
Different medium, same comparison label & 13 & 0.045 & 0.521 & 0.024 &
0.382 \\
Same medium, different label & 53 & 0.080 & 0.609 & 0.040 & 0.490 \\
Same medium, same comparison label & 12 & 0.121 & 0.643 & 0.045 &
0.402 \\
\end{longtable}

The same-medium/different-label MiniLM mean (0.609) exceeds the
cross-medium/same-label mean (0.521). This is consistent with a
source-context effect, but medium, date, editorial style, format,
document length and topic selection are confounded. It does not show
which medium is more authentic, that comparison labels are meaningless,
or that authors deliberately copied one another. No inferential
confidence intervals or p-values are assigned to these dependent pair
summaries.

\begin{figure}[p]
\centering
\includegraphics[width=1\textwidth,height=\textheight]{figures/source-label-contrast.pdf}
\caption{Dependent pair means by medium context and comparison label.
The category denominators are displayed; the study does not treat 300
pairs as 300 independent observations.}
\end{figure}

\begin{figure}[p]
\centering
\includegraphics[width=1\textwidth,height=\textheight]{figures/representation-maps.pdf}
\caption{Two lossy representations of the same 25 editions, with a
portable source key. The maps support source navigation; distances and
visual clusters do not validate doctrine or identity.}
\end{figure}

The lexical display is an uncentered two-component truncated SVD; the
neural display is a centered two-component PCA. Their displayed
explained-variance sums are approximately 6.30\% and 26.52\%,
respectively. Trustworthiness at five neighbors is 0.801 and 0.845. Both
coordinate systems are centered and scaled to unit root-mean-square
radius; orthogonal Procrustes rotates or reflects the neural display
only. The alignment RMS of 0.907 is not a latent-space equivalence
metric. Axes are not assigned metaphysical meanings.

\subsection{4.3 Stability and lexical
overlays}\label{stability-and-lexical-overlays}

Pair-rank Spearman between TF-IDF and MiniLM is 0.542. Mean top-five
neighbor Jaccard is 0.340. Full-versus-span overlap is 0.403 for TF-IDF
and 0.550 for MiniLM. Name masking produces higher neighbor stability
(0.924 lexical; 0.897 neural), but does not remove all source cues.
Equal-versus-content-weighted pooling changes at least one pair cosine
by as much as 0.093.

\begin{figure}[p]
\centering
\includegraphics[width=1\textwidth,height=\textheight]{figures/neighbor-stability.pdf}
\caption{Per-edition top-five neighbor stability across representation,
excerpt and masking choices. Set overlap measures modeling sensitivity,
not answer correctness.}
\end{figure}

\begin{longtable}[]{@{}
  >{\raggedright\arraybackslash}p{(\columnwidth - 8\tabcolsep) * \real{0.2000}}
  >{\raggedright\arraybackslash}p{(\columnwidth - 8\tabcolsep) * \real{0.2000}}
  >{\raggedright\arraybackslash}p{(\columnwidth - 8\tabcolsep) * \real{0.2000}}
  >{\raggedright\arraybackslash}p{(\columnwidth - 8\tabcolsep) * \real{0.2000}}
  >{\raggedright\arraybackslash}p{(\columnwidth - 8\tabcolsep) * \real{0.2000}}@{}}
\toprule\noalign{}
\begin{minipage}[b]{\linewidth}\raggedright
Representation
\end{minipage} & \begin{minipage}[b]{\linewidth}\raggedright
Silhouette
\end{minipage} & \begin{minipage}[b]{\linewidth}\raggedright
Medium ARI
\end{minipage} & \begin{minipage}[b]{\linewidth}\raggedright
Text-origin ARI
\end{minipage} & \begin{minipage}[b]{\linewidth}\raggedright
Comparison-label ARI
\end{minipage} \\
\midrule\noalign{}
\endhead
\bottomrule\noalign{}
\endlastfoot
Lexical & 0.036--0.037 & 0.351--0.443 & -0.072---0.034 & 0.029--0.113 \\
Semantic & 0.145--0.187 & 0.167--0.226 & 0.160--0.437 & 0.025--0.122 \\
\end{longtable}

The table shows ranges across the three fixed seeds. Low lexical
silhouettes and changing adjusted Rand indices caution against naming
clusters as stable civilizations or belief systems. Neural separation
also aligns with text origin in some runs. The lexical overlay averages
L2-normalized TF-IDF feature weights per medium with equal edition
weights, selecting the union of each medium's top six observed features,
capped at 30. It is an observed unigram/bigram overlay, not a neural
word-vector map, a prevalence estimate or a unique vocabulary test.

A separate synthetic control compares five constructed statements
outside the corpus. ``Contact will happen this year'' and its negated
version have MiniLM cosine approximately 0.830, exceeding the
approximately 0.542 score of the selected paraphrase. This small
demonstration illustrates why cosine cannot certify agreement. It is not
an accuracy benchmark or evidence that a corpus source contradicts
another.

\subsection{4.4 Contextual target-token
exploration}\label{contextual-target-token-exploration}

A separate add-on addresses literal word-vector overlays without
relabelling document embeddings or TF-IDF weights as word embeddings. It
selects 24 unigrams from 806 eligible candidates by mean fixed-baseline
TF-IDF weight, with document frequency at least three and all components
of exact medium and printed-label names excluded. The selection spans
409 source-term pairs in the same 25 editions. Up to two deterministic
ordinal occurrences per source and term yield 711 sampled contexts out
of 2,290 available occurrences. A single occurrence is used when only
one exists; the two-occurrence rule selects the first-third and
second-third ordinal positions rather than optimizing visible
separation.

For each occurrence, the same pinned MiniLM encoder processes up to 24
wordpieces on each side of the target. Final-layer hidden states for the
target's own subwords are averaged and normalized. Occurrences are
averaged equally within each source and normalized; source means are
then averaged equally for each term and normalized. This limits
domination by a source with many repeated occurrences, but medium
coverage still differs by term. These are exploratory target-token
states from a sentence-trained encoder, not validated word-sense
embeddings. Foundational work on contextual representation geometry
{[}R15{]} was inspected at abstract level; its BERT/ELMo/GPT-2 findings
are not presented as measurements of this MiniLM add-on.

One PCA fitted to the 24 term centroids transforms all 409 source-term
means in the same coordinate system. The two displayed components
explain 23.16\% of centroid variance, with trustworthiness at 5
neighbors of 0.8646. The first four selected terms---not, time, energy,
feel---supply the fixed source-overlay panels. Axes are shared within
this add-on only and are not aligned to the document maps. Fixed
three-cluster silhouettes range from 0.140 to 0.152; seed-pair adjusted
Rand indices range from 0.552 to 0.769. These are sensitivity
diagnostics, not an estimated number of meanings or civilizations.

\begin{figure}[p]
\centering
\includegraphics[width=1\textwidth,height=\textheight]{figures/contextual-token-overlay.pdf}
\caption{Selected target-token centroids and source-context overlays in
a shared PCA. Sentence-trained token states, two-occurrence sampling and
fixed clusters are exploratory; proximity is not common-answer agreement
or an ontology.}
\end{figure}

Frequent shared English vocabulary is overrepresented; sparse words,
untranslated material and unsampled uses remain outside this view.
High-cosine pairs include energy/energies (0.867), simply/just (0.778),
not/no (0.745). Morphology and ordinary language usage are competing
explanations for these neighborhoods, rather than evidence of a shared
channeling teaching. Token states encode spelling, local context and
source style as well as possible semantic associations. Nearness cannot
establish a shared proposition, agreement, contradiction, or a verified
speaker. Source character/token locators and their hashes are retained
privately, with no corpus quotations in the public figure. The original
document representations and TF-IDF term overlay remain unchanged.

\section{5. Source-card retrieval
experiment}\label{source-card-retrieval-experiment}

\subsection{5.1 Design and denominators}\label{design-and-denominators}

Four source-informed audit queries were authored independently of the
card author and frozen before that author saw them. They ask about a
proposed council's initiators and timing, preparedness through sharing
rather than hoarding, personal boundaries and external affirmation, and
rapid change versus stable internal awareness. This is useful separation
of roles, but the questions were still informed by source reading. They
are not a human-blinded external held-out test.

The raw index contains 276 nonoverlapping 180-whitespace-word chunks
covering all 47,280 words. Twenty-five cards cover 14 sources.
Answer-only cards contain 573 words; qualified cards contain 1,458.
Lexical vocabulary/IDF are fitted on raw source chunks only, and cards
and queries are transformed without refitting. A source receives its
maximum unit cosine, so document length and card multiplicity can affect
its opportunity to score highly. The neural condition uses the same
pinned MiniLM representation contract. The 13 navigation questions are
not included in the indexed cards.

Each query has one documented known-support source. Other retrieved
sources are \textbf{unjudged}, not automatically irrelevant.
Known-support Hit@k asks whether that one known source appears among the
first k sources. Reciprocal rank is zero if the source is missing and
otherwise the inverse of its rank; MRR averages four such values.
Coverage asks whether that support exists in an index at all. A
199-repetition source-ID shuffle is a descriptive negative control, not
a p-value, sampling confidence interval or estimate of general task
accuracy.

\begin{longtable}[]{@{}llllll@{}}
\toprule\noalign{}
Model / condition & MRR & Hit@1 & Hit@3 & Hit@5 & Support coverage \\
\midrule\noalign{}
\endhead
\bottomrule\noalign{}
\endlastfoot
TF-IDF / Raw chunks & 0.667 & 2/4 & 4/4 & 4/4 & 4/4 \\
TF-IDF / Answer cards & 0.750 & 3/4 & 3/4 & 3/4 & 3/4 \\
TF-IDF / Qualified cards & 0.625 & 2/4 & 3/4 & 3/4 & 3/4 \\
TF-IDF / Naive full-list RRF & 0.767 & 3/4 & 3/4 & 3/4 & 4/4 \\
TF-IDF / Top-5 RRF repair* & 0.800 & 3/4 & 3/4 & 4/4 & 4/4 \\
MiniLM / Raw chunks & 0.812 & 3/4 & 3/4 & 4/4 & 4/4 \\
MiniLM / Answer cards & 0.625 & 2/4 & 3/4 & 3/4 & 3/4 \\
MiniLM / Qualified cards & 0.562 & 2/4 & 2/4 & 3/4 & 3/4 \\
MiniLM / Naive full-list RRF & 0.766 & 3/4 & 3/4 & 3/4 & 4/4 \\
MiniLM / Top-5 RRF repair* & 0.786 & 3/4 & 3/4 & 3/4 & 4/4 \\
\end{longtable}

\emph{The top-five RRF repair is post-hoc. All fractions have four
queries as the denominator. Cards cover only three of their four known
support sources, imposing a 3/4 ceiling independently of ranking
quality.}

\begin{figure}[p]
\centering
\includegraphics[width=1\textwidth,height=\textheight]{figures/retrieval-audit.pdf}
\caption{Known-support retrieval for four fixed source-informed queries.
Naive full-list fusion is a diagnosed negative control; the top-five
repair was added after the first run. Rankings outside the known support
are unjudged.}
\end{figure}

\subsection{5.2 A structural coverage bias in naive
fusion}\label{a-structural-coverage-bias-in-naive-fusion}

For a source appearing in a ranked list, reciprocal-rank fusion
contributes \(1/(60+r)\), where \(r\) starts at one. In the naive
condition, all 25 raw sources and all 14 card-covered sources contribute
their full lists. Even the worst-positioned card-covered source receives

\[\frac{1}{60+25}+\frac{1}{60+14} \approx 0.02528,\]

whereas the best raw-only source receives at most

\[\frac{1}{60+1} \approx 0.01639.\]

Thus \textbf{every card-covered source outranks every raw-only source by
construction}, regardless of the query. This is a coverage artifact, not
an empirical discovery that synthesis improves relevance. An independent
critic derived the issue without viewing rankings, but the initial
experiment had already run. The correction therefore remains explicitly
post-hoc.

The repair fuses fixed top-five candidates from each list, then appends
the raw-ranked tail to retain coverage. It does not tune k over the four
queries or rewrite questions/cards. TF-IDF Hit@5 rises from 3/4 in naive
fusion to 4/4 in the repair; MiniLM remains 3/4. Raw MiniLM itself
reaches 4/4 at five with MRR 0.8125. These mixed, tiny-denominator
results do not select a universal winner or prove that qualified cards
improve retrieval. The operational lesson is to audit coverage and rank
mechanics before interpreting an apparent consolidation gain.

\section{6. Parallel speech-recognition
pilot}\label{parallel-speech-recognition-pilot}

\subsection{6.1 Feasibility and actual
performance}\label{feasibility-and-actual-performance}

The candidate is NVIDIA \texttt{parakeet-tdt-0.6b-v3}, revision
\nolinkurl{541d1f99c6b0c3cd0b11a95167540bb8edefd82b}. Its safetensors
artifact is 2,508,311,120 bytes; all downloaded model artifacts occupy
approximately 2.509 GB temporarily in RAM. The counted floating-state
elements are not presented as an independently verified architecture
parameter count. The choice is a strong contemporary open-model
candidate that runs in this environment, not a universal
state-of-the-art finding. Its joint technical report {[}R16{]},
inspected here at abstract and bibliographic level, announces Parakeet
multilingual recognition across 25 languages. That abstract's
1.7-million-hour training and ten-times-faster-than-Whisper claims
concern \textbf{Canary}, and are not transferred to this Parakeet run.

The pilot covers all four Council recordings plus four eligible
recordings from other contexts under a fixed selection rule. Two
nonoverlapping 30-second windows are centered near 25\% and 75\% of each
recording, selected before recognition. Six source contexts contribute
eight recordings and 16 paired windows: eight unique audio minutes and
32 model-window inferences. Both models receive the same canonical PCM
samples. Two workers each use one CPU thread on a shared host with a
four-core quota, 32 GiB memory limit and no available CUDA GPU.

\begin{longtable}[]{@{}lll@{}}
\toprule\noalign{}
Measured quantity & Whisper tiny.en & Parakeet 0.6b-v3 \\
\midrule\noalign{}
\endhead
\bottomrule\noalign{}
\endlastfoot
Concurrent workers / threads each & 2 / 1 & 2 / 1 \\
Outer-stage wall time, seconds & 12.48 & 62.40 \\
Observed inference overlap, seconds & 9.95 & 54.11 \\
Median 30-second-window inference, seconds & 1.248 & 6.856 \\
Median per-window real-time factor & 0.0416 & 0.2285 \\
Recognized whitespace words & 1015 & 913 \\
Conservative ASR memory accounting, GiB & 3.26 & 8.54 \\
\end{longtable}

The observed overlap proves that workers actually ran concurrently.
There is no matched serial execution on the same workload, so no
parallel speedup claim is made. Outer-stage time includes load and
orchestration; median real-time factor is a per-window measure. The
Parakeet run is slower in this pilot, although it remains faster than
real time per 30-second window. Shared host activity, runtime libraries,
decoding and model architectures differ; these timings do not isolate a
single cause or generalize to GPUs. Conservative memory accounting
includes worker RSS and staged artifacts and can overcount shared pages.

\subsection{6.2 Disagreement and timestamp
triage}\label{disagreement-and-timestamp-triage}

Symmetric normalized word-edit distance divides unit-cost token
Levenshtein distance by the longer output token sequence. Its median is
0.0706 and mean 0.1496; median distinct-token Jaccard is 0.8835. Two
windows differ in negation-token counts and zero in numeric-token
counts. These flags prioritize inspection; they do not establish
semantic contradiction or numerical correctness. Neither output is an
independent reference, so these metrics are \textbf{not word error rate
(WER), accuracy or hallucination rate}.

The retained tiny.en segment output overruns the 30-second window
endpoint in six segments across six windows, by at most one second. Raw
timestamps are preserved rather than silently repaired. Parakeet here
supplies window boundaries rather than comparable word-level alignment;
the pilot cannot rank alignment accuracy. Repeated output, omitted
speech, boundary artifacts and uncommon names require direct audio
review before substantive claim comparison. No such human listening
validation is claimed in this report.

A post-hoc repetition screen flags 1 baseline window and 0 Parakeet
windows at a threshold of at least ten identical consecutive tokens. In
the flagged Ra window, the baseline's longest identical-token run
contains 112 tokens. All 16 windows remain in the reported summaries.
This is an output-quality triage flag; without independent listening it
is not a confirmed hallucination, a diagnosis of the underlying audio,
or proof that the other model's output is correct.

The eight-recording recognition pilot and the 25-edition representation
frame are separate selections. For example, a pilot Ra window is not a
Ra edition in the main 25-source knowledge bank. Results from one
denominator must not silently be transferred to the other.

\section{7. Full-recording paired-edition
sensitivity}\label{full-recording-paired-edition-sensitivity}

The separately completed run covers \textbf{4 full recordings, 178.99
audio minutes and 359 contiguous nonoverlapping windows}, producing
26,319 recognized whitespace words. Exact canonical 16 kHz mono sample
bindings, start/end offsets and final partial tails are checked.
Complete sample coverage does not establish complete or accurate
linguistic recognition. Independent windows can split an utterance or
lose discourse context; output is concatenated in sample order without
inferred speaker identities or silently repaired wording. All 2
empty-text chunks are retained as empty outcomes rather than omitted
from coverage.

The measured two-worker stage takes 21.23 minutes, with one CPU thread
per worker and 20.88 minutes of observed inference overlap. Median chunk
real-time factor is 0.2344; conservative memory accounting is 9.33 GiB.
The stage includes hash checks, model loading and inference, and
excludes download, preparation and final projection. It is a shared-host
measurement without a matched serial speedup comparison.

For a controlled sensitivity analysis, four old recognition editions are
counterfactually replaced within the 25-source frame; 21 editions remain
unchanged. The published original fit is preserved. The new 25-source
counterfactual has 47,635 words versus 47,280; the paired subset changes
from 25,964 to 26,319. More words are not evidence of better
recognition.

The lexical transform freezes the original phrase mask, vocabulary and
IDF; new terms outside that vocabulary do not become new features. The
semantic transform uses the same pinned MiniLM and complete
chunk-pooling contract. This is a recognition-pipeline sensitivity
analysis: model, context/VAD handling, segmentation, decoding and
punctuation differ, so a change cannot be attributed solely to model
weights.

\begin{longtable}[]{@{}ll@{}}
\toprule\noalign{}
Observed sensitivity & Value \\
\midrule\noalign{}
\endhead
\bottomrule\noalign{}
\endlastfoot
Mean old/new paired TF-IDF cosine & 0.9616 \\
Mean old/new paired MiniLM cosine & 0.9880 \\
All-source mean top-five Jaccard, TF-IDF & 0.9600 \\
All-source mean top-five Jaccard, MiniLM & 0.8933 \\
Largest pair-cosine change, TF-IDF & 0.0109 \\
Largest pair-cosine change, MiniLM & 0.0259 \\
Existing cards requiring edition-specific re-reading & 9 \\
Cards automatically rewritten & 0 \\
\end{longtable}

\begin{figure}[p]
\centering
\includegraphics[width=1\textwidth,height=\textheight]{figures/asr-edition-sensitivity.pdf}
\caption{Four same-audio pairs under old and new recognition pipelines.
Word counts, representation cosines and normalized machine-text
disagreement describe sensitivity; neither edition is ground truth.}
\end{figure}

All existing cards remain bound to their old source wording and
locators. New editions require independent re-reading and re-annotation
before inheriting any claim or relation. Similarity changes do not
automatically alter a compatibility or contradiction label. Four
deliberately selected Council recordings cannot establish performance
across all channels, languages or recording conditions.

Independent model-agent review checked every retained input/output join
and all 359 window-duration bindings. A separate scalar
dynamic-programming implementation reproduced full token-edit distances
of 104, 145, 108 and 1,602. An independent log-TF/IDF reconstruction
reproduced the baseline and counterfactual lexical matrices; arithmetic
on saved neural-vector witnesses reproduced paired cosines, matrices and
neighbor summaries. These checks did not repeat neural-model inference,
establish speaker identity, or supply a listening-verified reference.
Detailed acceptance scopes are retained in the project's
independent-review record.

\section{8. Cross-field connections and falsifiable next
experiments}\label{cross-field-connections-and-falsifiable-next-experiments}

The useful cross-field transfers are methodological, and each requires a
stated failure condition.

\textbf{Information retrieval and archive coverage.} Card compression
can change both relevance and opportunity to be retrieved. A larger
preregistered query set should stratify supports that have and lack
cards, hold query families out of tuning, report known-support coverage
separately, and audit multiple plausible supports rather than treating
all unlisted sources as irrelevant. A claimed gain fails if it
disappears under equal source coverage or depends on repeatedly tuning
fusion to the same questions. The present four-query diagnostic
identifies a structural failure and supplies no population accuracy
estimate.

\textbf{Oral history, discourse analysis and claim scope.} A statement's
practical meaning depends on audience, speaker attribution, date,
referent, condition and modality. A blinded annotation exercise can
present source spans, qualified cards and answer-only cards to
independent reviewers and count which obligatory qualifiers survive. The
denominator should be prespecified qualifier obligations, with
disagreement and unresolved cases retained. A compression method fails
if it systematically turns possibilities into certainties or collective
metaphors into dated factual predictions. Current cards are source-bound
but have not undergone a human inter-annotator study.

\textbf{Stylometry and measurement theory.} Medium and text-origin
signals can masquerade as civilization/topic similarity. A balanced
future design should sample several matched questions per medium and
label, preserve publication and editing metadata, and evaluate held-out
medium or publisher conditions. A civilization-level interpretation
weakens if separation vanishes after matching genre, length and topic or
if medium prediction dominates. Current dependent-pair contrasts are
descriptive; the available sample does not permit independent causal
separation of those factors.

\textbf{Memory consolidation and regression testing.} Immutable sources,
compressed cards and a retrieval index form three inspectable layers.
Repeated updates can retain a fixed old-question set alongside newly
introduced questions, reporting source support, qualifier retention,
unsupported synthesis and index coverage after each update. A
consolidation proposal fails if gains on new questions hide losses on
old questions or erase provenance. This is an external-memory regression
design, not evidence that a language model has acquired new parametric
knowledge.

\textbf{Speech recognition and error propagation.} A future
listening-verified reference should stratify quiet speech, noise,
repetition, unfamiliar names, negation and numerical claims. Word-level
accuracy, named-entity correctness and source-card stability should be
assessed separately, using the same audio and explicit pipeline
settings. Current machine-to-machine disagreement identifies review
targets but cannot show which model is right. A higher-capacity model
fails the substantive task if it produces fluent but unsupported names
or reverses conditions even when its broad topic embedding remains
similar.

\textbf{Network epistemology and graph semantics.} Typed provenance
paths can distinguish ``source supports card,'' ``analyst compares
propositions'' and ``hypothesis connects methods.'' A later
graph-retrieval experiment can compare equal-budget lexical, neural and
graph retrieval against documented supports. A graph method fails if
apparent gains arise from duplicated source edges or label leakage. High
degree and dense communities are not independent corroboration. GraphRAG
and hippocampal-memory papers motivate tests, not shortcuts to
evidential authority.

\textbf{Forecast auditing.} Date-bound prediction/outcome matching
remains proposed. It requires the original publication timestamp, stable
edition text, exact forecast scope, outcome criteria selected before
outcome lookup, and listening verification where the decisive date comes
from ASR. Retrospective thematic similarity is insufficient. The present
study does not score prophecy, supernatural causation, health outcomes
or future events.

\section{9. Limitations, AI roles and
reproducibility}\label{limitations-ai-roles-and-reproducibility}

This study is constrained by purposive source access, few mediums,
uneven labels and topics, variable edition lengths and mixed text
origins. Publisher text is not an independent word-aligned audio
reference. Embeddings encode training-distribution biases; numerical
claims and negation require direct examination. Lossy maps, selected
lexical features and fixed cluster counts are exploratory. The pair
observations are dependent. Four retrieval queries and four
full-recognition source pairs cannot support general superiority claims.
No new weight training, human-blinded accuracy study, verified channeled
identity or causal spiritual mechanism is established.

AI agents performed literature inspection, source reading and bounded
paraphrase annotation, representation design, code implementation, ASR
orchestration, report writing and independent criticism in distinct
roles. Independent model-agent review is not independent human
replication or inter-annotator agreement. The source-reading and card
roles were separated from the four audit-query authorship role before
the card author saw the queries. An independent critic identified
qualifier issues and the naive-fusion bias; the latter remained post-hoc
because the first run already existed. Runtime and script/model/data
hashes bind executed outputs. Maintainer authorship and human validation
remain TBD.

The report is additive. Earlier source editions, numeric fits, papers
and raw evidence are retained. Private full text and audio are excluded
from the public source archive. Public materials contain the analysis
code, literature register, bounded card paraphrases, question database,
typed graph, figures, this report and release hashes. Access
restrictions and source licences continue to apply to raw materials. The
report source archive is self-contained for typesetting, not a
redistribution of the research corpus or downloaded model weights.

To reproduce the PDF from the source archive, run
\texttt{pdflatex\ -interaction=nonstopmode\ -halt-on-error\ main.tex}
twice in a directory containing its \texttt{figures/} folder. To
regenerate the report's selected numbers, run \texttt{build.py} with the
accepted analysis, knowledge, retrieval, ASR and literature JSON inputs
described in the accompanying README. Full-recording results require
both a completed manifest and its hash-bound sensitivity artifact.
Missing inputs produce an explicitly pending chapter, never inferred
completion. The build verifies selected figure hashes and pins the main
numerical and retrieval releases.

\subsection{9.1 Input bindings}\label{input-bindings}

\begin{longtable}[]{@{}lll@{}}
\toprule\noalign{}
Accepted input role & SHA-256 prefix & Bytes \\
\midrule\noalign{}
\endhead
\bottomrule\noalign{}
\endlastfoot
analysis & e1407aa84f8eb333 & 487,845 \\
knowledge & 9ef331231f1efdb5 & 181,571 \\
retrieval & 0df6a4967ed08fd1 & 251,376 \\
asr & 11a4242e924fa250 & 42,359 \\
literature & 639e8cb68aed09c6 & 27,061 \\
asr\_full & 3991b8649962b522 & 265,080 \\
asr\_full\_public & fdfdd4d7fdeec4c7 & 10,851 \\
sensitivity & 02aacfb8ceef15df & 81,350 \\
contextual\_terms & 5d3d509f6a509270 & 159,262 \\
additional\_references & 3b0a9be8df956874 & 1,852 \\
\end{longtable}

Full SHA-256 values and output/figure bindings are in
\texttt{manifest.json}. Prefixes above are display abbreviations, not
alternate hashes.

\section{References}\label{references}

\textbf{R1. Ali Behrouz; Farnoosh Hashemi; Adel Javanmard; Vahab
Mirrokni.} Language Models Need Sleep: Learning to Self-Modify and
Consolidate Memories. arXiv:2606.03979v2; revised 2026-07-10. Primary
source: \url{https://arxiv.org/html/2606.03979v2}. Access depth: full
text; retrieved 2026-11-08. Inspection details and source hashes are in
the literature register.

\textbf{R2. Ali Behrouz; Vahab Mirrokni.} The sleep framework: A new way
for LLMs to retain knowledge in continual learning. Google Research
blog, retrieved 2026-11-08. Primary source:
\url{https://research.google/blog/the-sleep-framework-a-new-way-for-llms-to-retain-knowledge-in-continual-learning/}.
Access depth: full text; retrieved 2026-11-08. Inspection details and
source hashes are in the literature register.

\textbf{R3. Kevin Lin; Charlie Snell; Yu Wang; Charles Packer; Sarah
Wooders; Ion Stoica; Joseph E. Gonzalez.} Sleep-time Compute: Beyond
Inference Scaling at Test-time. arXiv:2504.13171v1. Primary source:
\url{https://arxiv.org/html/2504.13171v1}. Access depth: full text;
retrieved 2026-11-08. Inspection details and source hashes are in the
literature register.

\textbf{R4. Ali Behrouz; Meisam Razaviyayn; Peilin Zhong; Vahab
Mirrokni.} Nested Learning: The Illusion of Deep Learning Architectures.
arXiv:2512.24695v1; Google blog dated 2025-11-07. Primary source:
\url{https://arxiv.org/abs/2512.24695}. Access depth: abstract or
excerpt; retrieved 2026-11-08. Inspection details and source hashes are
in the literature register.

\textbf{R5. Darren Edge; Ha Trinh; Newman Cheng; Joshua Bradley; Alex
Chao; Apurva Mody; Steven Truitt; Dasha Metropolitansky; Robert Osazuwa
Ness; Jonathan Larson.} From Local to Global: A Graph RAG Approach to
Query-Focused Summarization. arXiv:2404.16130v2; revised 2025-02-19.
Primary source: \url{https://arxiv.org/html/2404.16130v2}. Access depth:
full text; retrieved 2026-11-08. Inspection details and source hashes
are in the literature register.

\textbf{R6. Bernal Jiménez Gutiérrez; Yiheng Shu; Yu Gu; Michihiro
Yasunaga; Yu Su.} HippoRAG: Neurobiologically Inspired Long-Term Memory
for Large Language Models. arXiv:2405.14831v3; revised 2025-01-14.
Primary source: \url{https://arxiv.org/html/2405.14831v3}. Access depth:
full text; retrieved 2026-11-08. Inspection details and source hashes
are in the literature register.

\textbf{R7. Bernal Jiménez Gutiérrez; Yiheng Shu; Weijian Qi; Sizhe
Zhou; Yu Su.} From RAG to Memory: Non-Parametric Continual Learning for
Large Language Models. arXiv:2502.14802v2; revised 2025-06-19. Primary
source: \url{https://arxiv.org/abs/2502.14802v2}. Access depth: abstract
or excerpt; retrieved 2026-11-08. Inspection details and source hashes
are in the literature register.

\textbf{R8. Parth Sarthi; Salman Abdullah; Aditi Tuli; Shubh Khanna;
Anna Goldie; Christopher D. Manning.} RAPTOR: Recursive Abstractive
Processing for Tree-Organized Retrieval. arXiv:2401.18059v1. Primary
source: \url{https://arxiv.org/abs/2401.18059}. Access depth: abstract
or excerpt; retrieved 2026-11-08. Inspection details and source hashes
are in the literature register.

\textbf{R9. Kenneth Enevoldsen; Isaac Chung; Imene Kerboua; et al..}
MMTEB: Massive Multilingual Text Embedding Benchmark.
arXiv:2502.13595v4; revised 2025-11-13. Primary source:
\url{https://arxiv.org/html/2502.13595v4}. Access depth: full text;
retrieved 2026-11-08. Inspection details and source hashes are in the
literature register.

\textbf{R10. Nils Reimers; Iryna Gurevych.} Sentence-BERT: Sentence
Embeddings using Siamese BERT-Networks. arXiv:1908.10084v1. Primary
source: \url{https://arxiv.org/abs/1908.10084}. Access depth: abstract
or excerpt; retrieved 2026-11-08. Inspection details and source hashes
are in the literature register.

\textbf{R11. Marie-Catherine de Marneffe; Anna N. Rafferty; Christopher
D. Manning.} Finding Contradictions in Text. ACL-08: HLT,
pp.~1039--1047. Primary source:
\url{https://aclanthology.org/P08-1118/}. Access depth: full text;
retrieved 2026-11-08. Inspection details and source hashes are in the
literature register.

\textbf{R12. Alec Radford; Jong Wook Kim; Tao Xu; Greg Brockman;
Christine McLeavey; Ilya Sutskever.} Robust Speech Recognition via
Large-Scale Weak Supervision. arXiv:2212.04356v1. Primary source:
\url{https://arxiv.org/abs/2212.04356}. Access depth: abstract or
excerpt; retrieved 2026-11-08. Inspection details and source hashes are
in the literature register.

\textbf{R13. Timnit Gebru; Jamie Morgenstern; Briana Vecchione; Jennifer
Wortman Vaughan; Hanna Wallach; Hal Daumé III; Kate Crawford.}
Datasheets for Datasets. arXiv:1803.09010v8; revised 2021-12-01. Primary
source: \url{https://arxiv.org/abs/1803.09010v8}. Access depth: abstract
or excerpt; retrieved 2026-11-08. Inspection details and source hashes
are in the literature register.

\textbf{R14. Timothy Lebo; Satya Sahoo; Deborah McGuinness.} PROV-O: The
PROV Ontology. W3C Recommendation, 2013-04-30. Primary source:
\url{https://www.w3.org/TR/prov-o/}. Access depth: full text; retrieved
2026-11-08. Inspection details and source hashes are in the literature
register.

\textbf{R15. Kawin Ethayarajh.} How Contextual are Contextualized Word
Representations? Comparing the Geometry of BERT, ELMo, and GPT-2
Embeddings. EMNLP-IJCNLP 2019, pp.~55--65, DOI 10.18653/v1/D19-1006.
Primary source: \url{https://aclanthology.org/D19-1006/}. Access depth:
abstract and bibliographic metadata only; retrieved 2026-11-08.
Foundational caution that contextual representations have
model/layer-dependent geometry; the paper does not validate the present
MiniLM token means or clusters.

\textbf{R16. Monica Sekoyan; Nithin Rao Koluguri; Nune Tadevosyan; Piotr
Zelasko; Travis Bartley; Nikolay Karpov; Jagadeesh Balam; Boris
Ginsburg.} Canary-1B-v2 \& Parakeet-TDT-0.6B-v3: Efficient and
High-Performance Models for Multilingual ASR and AST.
arXiv:2509.14128v2, revised 25 September 2025. Primary source:
\url{https://arxiv.org/abs/2509.14128v2}. Access depth: abstract and
bibliographic metadata only; retrieved 2026-11-08. The abstract
announces Parakeet multilingual ASR across 25 languages; its 1.7M-hour
and ten-times-Whisper performance claims concern Canary, not this
Parakeet experiment.

\end{document}
